% Background and Motivation: trace structure and matched physical KV costs.
\section{Background and Motivation}
\label{sec:background-motivation}

\subsection{Background}
\label{sec:motivation-background}

\paragraph{Multi-agent workflows.}
Coding assistants and research systems increasingly organize LLM calls
into multi-agent workflows~\cite{wu2023autogen,anthropic2025multiagent}.
An agent may have a predefined role, such as planner or reviewer, or be
created for a particular subtask. It can issue multiple LLM requests,
invoke tools, and incorporate their results before producing an output.
Sequential dependencies, parallel branches, and joins connect these
agents. A request, an agent, and the complete user workflow therefore
have different completion conditions.

\paragraph{Prefill and reusable state.}
Prefill processes a request's input tokens and produces per-layer key and
value (KV) tensors. Subsequent computation can reuse this state instead
of processing the same prefix again. Agent histories often grow as later
requests include earlier reasoning, code, and tool outputs. Reusing their
KV saves computation, but retaining every growing history consumes memory.
Hierarchical caches extend GPU memory with host DRAM and larger storage
tiers~\cite{qin2025mooncake,gao2024cachedattention}.

\paragraph{Failures at serving scale.}
A GPU or worker failure can interrupt an entire model replica, removing
both serving capacity and access to its active KV cache~\cite{cao2026lumen}.
Surviving KV copies allow recovery without reconstructing the complete
history; otherwise, retained tokens must be processed again.
KV checkpointing and load-aware recovery address these
costs~\cite{strati2024dejavu,jayakody2026ghostserve,cao2026lumen}.
We consider fail-stop failures with a compatible surviving replica,
retained token history and control state, and backup storage outside the
failed domain. Within this scope, we ask which histories deserve scarce
fast-tier capacity when their recovery affects a dependent workflow.

\begin{figure*}[t]
  \centering
  \captionsetup[subfigure]{font=normalsize,labelfont=normalfont,
    textfont=normalfont,justification=centering,singlelinecheck=false}
  \begin{subfigure}[t]{2.48in}
    \vspace{0pt}
    \begin{minipage}[c][1.65in][c]{\linewidth}
      \centering
      \includegraphics[width=\linewidth]{figures/workflow.pdf}
    \end{minipage}
    \caption{Four-branch workflow}
    \label{fig:motivation-workflow}
  \end{subfigure}\hfill
  \begin{subfigure}[t]{1.86in}
    \vspace{0pt}
    \begin{minipage}[c][1.65in][c]{\linewidth}
      \centering
      \includegraphics[width=\linewidth]{figures/workflow-gaps.pdf}
    \end{minipage}
    \caption{Workflow-maximum gaps}
    \label{fig:motivation-branch-gap}
  \end{subfigure}\hfill
  \begin{subfigure}[t]{2.52in}
    \vspace{0pt}
    \begin{minipage}[c][1.65in][c]{\linewidth}
      \centering
      \includegraphics[width=\linewidth]{figures/recovery-cost.pdf}
    \end{minipage}
    \caption{Context and Prefill cost}
    \label{fig:motivation-recompute}
  \end{subfigure}
  \caption{Workflow structure and the cost of reusable KV state.
    (a) Four trace-derived branches with schematic Main stages A0/A5.
    Bars preserve timing; dashed tails assume an all-branch join.
    (b) ECDF of maximum branch-duration gaps per workflow: 121 workflows
    with temporally inferred fan-in; these are not failure-time slack.
    (c) Measured Qwen2.5-7B-Instruct-1M Prefill on two A100 PCIe 40\,GB GPUs
    (TP2, BF16, batch one), with sparse attention above 32K.
    Reuse starts from $(C-1)$ cached tokens and computes one;
    DRAM/shared-storage I/O is serialized with computation.
    HBM/DRAM: median and min--max of three timings;
    shared storage/full prefill: one timing per length.
    Protocol: Appendix~\ref{sec:motivation-protocol}.
    K~=~1,024; M~=~1,000,000.}

  \label{fig:motivation}
  \Description{Left: a four-branch workflow with A2 completing last at
    151.6 seconds. Middle: an ECDF of workflow-maximum branch-duration
    differences with a median of 138 seconds and a 90th percentile of
    1173 seconds. Right: four measured Prefill curves for Qwen2.5-7B-Instruct-1M
    from 8K to 1,010,000 tokens: blue circles for HBM reuse, orange squares
    for serialized DRAM reuse, purple diamonds for serialized shared-storage
    reuse, and red triangles for full prefill. Both axes are logarithmic.}
\end{figure*}

\subsection{Motivation}

\paragraph{Workflow completion requires different protection priorities.}
\label{sec:motivation-workflow}
We use workflow completion time (WCT), from task release to completion of
all required results, as the application-level objective.
Figure~\ref{fig:motivation}(\subref{fig:motivation-workflow}) illustrates
why request latency alone is insufficient. A2 finishes last, 151.6\,s
after the fork, whereas A1 finishes at 60.4\,s. If the following Main stage
requires all four outputs, accelerating A1 alone does not advance the
join, while delaying A2 postpones downstream work. A branch can absorb
some additional delay without postponing the required output; we call
this tolerance \emph{slack}.

The public SemiAnalysis CC/Weka trace contains 393 filtered Claude Code
sessions~\cite{semianalysis2026cctraces}. We retain overlapping subagent
groups whose completed branches are followed by a Main request after the
last branch finishes; Main may also execute while the branches run.
This temporal fan-in proxy retains 286 groups in 121 workflows.
For each group we take the longest minus the shortest branch duration,
then retain each workflow's maximum gap.
Figure~\ref{fig:motivation}(\subref{fig:motivation-branch-gap}) gives a
median of 138\,s and a 90th percentile of 1,173\,s. These conditional
statistics establish substantial branch-duration heterogeneity. They do
not establish result-consumption edges, which the trace does not record,
or the slack distribution at a fault. The four-branch example illustrates
one retained group, not a representative sample of workflow maxima.

\paragraph{Long contexts make fast recovery state expensive.}
\label{sec:motivation-long-context}
Losing an agent's KV can turn a brief interruption into substantial
reconstruction work. Figure~\ref{fig:motivation}(\subref{fig:motivation-recompute})
compares four paths using Qwen2.5-7B-Instruct-1M~\cite{qwen2025qwen25oneM}
on two A100 PCIe 40\,GB GPUs, with TP2, BF16 and batch one.
All paths use the same SGLang 0.5.10rc0 model execution and
MInference vertical/slash attention~\cite{jiang2024minference} with dual chunk
attention (DCA) position handling; attention is sparse above 32K.
For total context $C=M+1$, full prefill starts from empty KV, while reuse
starts from the same $M$-token history in HBM, DRAM, or shared storage and
computes one token. Timing ends with all $C$ tokens' KV ready, before
vocabulary projection. HBM provides a resident-state reference; the
DRAM and shared-storage paths serialize restoration and computation.

At 1,010,000 tokens, full prefill takes 571.21\,s despite sparse attention.
HBM, DRAM and shared-storage reuse take 1.305, 2.663 and 10.156\,s,
respectively. A reusable history avoids most reconstruction work, but its
storage tier still changes readiness time: shared storage adds 7.49\,s
over DRAM in this run. The corresponding history occupies 57.92\,GB
across the two ranks. Keeping many such histories in DRAM can therefore
be expensive, while moving them to storage introduces additional delay.
These measurements characterize one fixed synthetic input and serialized
paths, with the timing counts stated in the caption. They establish the
cost difference in this configuration; storage-cache state and overlapping
I/O can change its magnitude. Appendix~\ref{sec:motivation-protocol}
gives the measurement and validation details.

\begin{figure}[t]
  \centering
  \includegraphics[width=\linewidth]{figures/recovery-exposure.pdf}
  \caption{Measured Prefill costs for
    Figure~\ref{fig:motivation}(\subref{fig:motivation-workflow}), using
    Figure~\ref{fig:motivation}(\subref{fig:motivation-recompute})'s configuration
    at trace input-size proxies plus one token.
    Bars/whiskers show median/min--max over three timings after warmup.
    Hatching denotes absorbed delay; solid fill is the calculated WCT
    increase under an all-branch join with other execution fixed.
    Each case is a single-branch counterfactual, not a fault-injected
    WCT measurement.}
  \label{fig:motivation-recovery-exposure}
  \Description{A1 through A4 have measured median Prefill costs of
    33.2, 32.6, 35.4 and 39.9 seconds.
    Applying each cost separately as an added branch pause gives WCT
    increases of 0.0, 32.6, 16.3 and 18.0 seconds.
    Hatching identifies the absorbed portion; whiskers show the minimum
    and maximum of three timings, not confidence intervals.}
\end{figure}

\paragraph{Failure position determines the workflow benefit of protection.}
Reconstruction cost and workflow position must be considered together.
Figure~\ref{fig:motivation-recovery-exposure} uses the four branches in
Figure~\ref{fig:motivation}(\subref{fig:motivation-workflow}). We measure
full prefill at each active request's input-size proxy, approximately
101K--114K tokens in decimal units, using the configuration above.
The proxies come from trace hash-block counts; the measured input is
synthetic, rather than the original Claude prompts.
At 29.5\,s after the fork, we insert each measured cost $R_i$ as an
additional pause on branch $i$, separately, while holding other execution
and queues fixed. Under
the assumed all-branch join, the trace completion times give offline
slack $\sigma_i$, and the resulting WCT increment is
$\max(0,R_i-\sigma_i)$.

A1's 33.2\,s cost fits within its 91.2\,s slack and adds no WCT
in this construction. A2 takes 32.6\,s but has zero slack,
exposing its entire cost. A3 and A4 expose 16.3 and
18.0\,s, respectively. Thus, reconstruction latency alone does
not determine the benefit of protection. This example combines physical
reconstruction timings with a conditional dependency calculation; it
does not measure failures in the original workflow or establish the
frequency of such outcomes. Common-worker failures and shared recovery
resources require joint analysis, developed in the next section.

\label{sec:motivation-goals}
These results motivate AegisKV's central objective: allocate finite
fast-tier capacity and backing storage according to their effect on
workflow completion. Retaining a usable copy avoids reconstruction;
keeping KV closer to execution can additionally avoid restore waits.
Neither saving has the same WCT value for every branch. Protection must also account for normal reuse
and the capacity and I/O it consumes. The design therefore separates the
value of fast-tier retention from the value of keeping any recoverable
copy, and evaluates both in the context of the workflow.
